\begin{result}{F.3.13}{The integral automorphism model of an admissibly generated finite stage}{Proved for the specified exact category}
Let $R$ be noetherian, let $\mathcal C$ be a small rigid symmetric $R$-linear exact tensor category, and let
$\omega:\mathcal C\to\operatorname{Proj}_f(R)$ be an exact faithful tensor functor. Fix an object $M$ from which every object of $\mathcal C$ is obtained, up to isomorphism, by finitely many direct sums, tensor products, duals, and admissible subobjects and quotients. For an $R$-algebra $A$, put $\omega_A(X)=\omega(X)\otimes_R A$. Restriction to $M$ represents the functor
\[
 A\longmapsto\operatorname{Aut}^{\otimes}(\omega_A)
\]
by a closed subgroup scheme of $\operatorname{GL}_R(\omega M)$, hence by an affine group scheme of finite presentation over $R$.

For $R=\mathcal O_{F,S}$, with $F$ a number field and $S$ a finite set of finite primes, the isomorphism classes of fpqc-local tensor forms of this particular $\omega$ are consequently finite, by the explicit torsor equivalence and Dedekind-ring approximation in F.3.12. No smoothness, flatness, connectedness or reductivity of the resulting automorphism scheme is needed.
\end{result}
\begin{proof}
Write $V=\omega M$ and $C=\mathcal O(\operatorname{GL}_R(V))$. The universal automorphism $u$ and its inverse have coefficients in $C$. For each object choose one finite expression by the indicated operations, together with the chosen admissible sequences and final isomorphism. Tensor, dual and sum operations define matrices polynomial in the coefficients of $u,u^{-1}$.

At an admissible subobject $i:N\hookrightarrow P$, exactness gives
\[
 0\longrightarrow\omega N\xrightarrow{\omega i}\omega P\longrightarrow Q\longrightarrow0
\]
with $Q$ finite projective. This sequence splits as $R$-modules. Choose a retraction $r:\omega P\to\omega N$, and impose the polynomial equations
\[
 (1-(\omega i)r)\,u_P(\omega i)=0,
 \qquad
 (1-(\omega i)r)\,u_P^{-1}(\omega i)=0.
\]
Their meaning over every $A$ is preservation of $\omega_A N$ by $u_P$ and $u_P^{-1}$. The induced mutually inverse matrices on $\omega_A N$ are $r u_P(\omega i)$ and $r u_P^{-1}(\omega i)$. The split injection remains injective after every base change, so these are the unique matrices making naturality for $i$ hold. At an admissible quotient $q:P\twoheadrightarrow N$, its projective kernel and projective quotient similarly give a split sequence. Preservation of the kernel by $u_P,u_P^{-1}$ gives a unique induced automorphism; for a chosen section $s$ its matrix is $(\omega q)u_Ps$. These formulas handle any finite expression inductively, including successive subquotients. Finite projective modules may be written as images of idempotent matrices in finite free modules, so all displayed preservation conditions are ordinary coefficient equations even without a global basis.

Now impose, for every morphism $a:X\to Y$ of $\mathcal C$, the equations
\[
 u_Y\omega(a)=\omega(a)u_X,
\]
and impose the tensor, unit and dual compatibility equations. In particular these impose compatibility between the independently chosen expressions for $X\otimes Y$ and the tensor product of the chosen expressions for $X,Y$. Let $I\subset C$ be the ideal generated by all these equations and the subquotient preservation equations. The construction gives, for every $R$-algebra $A$, mutually inverse maps between $\operatorname{Hom}_{R\text{-alg}}(C/I,A)$ and tensor automorphisms of $\omega_A$: a point supplies the constructed components; an automorphism supplies its component at $M$; naturality and split exactness force all its other components. Thus the assertion is scheme-theoretic and includes algebras with nilpotents and torsion. It is not merely a calculation on field-valued points.

The algebra $C$ is finitely presented over noetherian $R$, hence noetherian. Its ideal $I$ is finitely generated. Therefore $\operatorname{Spec}(C/I)$ is of finite presentation over $R$. Composition and inversion of the natural automorphisms restrict those of $\operatorname{GL}(V)$, giving the claimed closed subgroup scheme. Finally apply F.3.12 to this actual automorphism scheme and the actual map $\omega'\mapsto\operatorname{Isom}^{\otimes}(\omega,\omega')$.
\end{proof}

\paragraph{Historical boundary.}
This construction removes the extra representability hypothesis for the exact finite stages just defined. The supplied 15:213 uses $\mathbb M_\alpha$ and $T_\alpha$, and 15:223 calls $G_M$ an algebraic group over $\mathbb Z$. Identifying those historical stages with this admissibly generated integral category still requires their definition; finite generation of the rational category alone does not supply the split integral presentations used in this proof. The unrestricted proalgebraic tensor category in F.3.12 has explicitly infinitely many forms.
